\section{Discussion}
\label{sec:discussion}

\paragraph{What the results say about defense placement.}
Our results indicate that the marginal value of a defense layer depends on where it sits in the execution lifecycle. Prompt-hardening and information-flow defenses act before or during a single episode and cannot learn from episodes that have already completed; \method{} acts after completion and converts residual risk---the attacks that slipped past everything else---into an asset. The paired replay design makes this value proposition measurable: because the same trajectory is re-executed with and without the injected experience under a fixed controller, the ASR difference is attributable to the experience content rather than to case mix, decoding stochasticity, or prompt-length effects. The ablation pattern further localizes the source of improvement: replacing the graph-localized evidence with raw trajectory text or with a generic safety instruction weakens the effect, which supports our claim that \emph{what} is injected matters as much as \emph{whether} anything is injected.

\paragraph{Label-free operation as a deployment property.}
A defining constraint of \method{} is that no ground-truth label reaches the LLM at any stage: the detector trains on graph-level labels offline, the scorer inherits weak MIL supervision, the judge is instructed to reject benign and false-positive cases without ever seeing labels, and the replay judge scores outcomes it did not produce. This matters for practice because production deployments cannot afford labeling of every suspicious trace, and because any label leak would trivially inflate experience quality. The cost of this constraint is visible in the detection numbers: at a 15\% abnormal base rate, PR-AUC is moderate and the threshold favors precision, so some attacks pass unflagged and receive no experience. The pipeline therefore behaves as a \emph{partial} postmortem layer whose coverage grows with detector recall---an explicit, quantifiable trade-off rather than a hidden failure mode.

\paragraph{Judge dependence and its mitigation.}
Both the synthesis judge and the attack-success judge are LLMs, and the reliability literature documents verdict instability even under pinned decoding~\cite{zhou2026style,ayyagari2026pinned}. We mitigate this in three ways: model-family separation between agent, synthesizer, and judge; spot-checking judge verdicts against the benchmark's recorded compromised executions; and reporting agreement statistics rather than treating the judge as an oracle. Residual risks remain. A judge that systematically under-detects compromise in the presence of well-formed final answers would deflate ASR for both conditions and compress our measured improvement; the trajectory-level audits partially control for this because they anchor verdicts to recorded outcomes, but a fully human-audited subset is the proper instrument and is left for future work.

\paragraph{Scope and limitations.}
Three limitations bound our claims. First, all experiments run on \tool{}; the pipeline is model- and dataset-agnostic in construction---the trajectory-graph schema, the detector interface, and the replay protocol are portable, as evidenced by the bundled ATBench configuration---but cross-dataset transfer of experiences is untested here because the ATBench task corpus was unavailable in usable form during this study. Second, experiences are synthesized from the same distribution they defend against; our within-case design maximizes measurement sensitivity but optimistic-bounds the cross-case transfer we report as a secondary analysis. Third, the replay agent observes a sanitized environment with recorded tool responses; adaptive adversaries who mutate payloads between the baseline and experience runs could erode the measured protection, and the redaction layer, while covering credential-shaped strings, cannot guarantee the absence of all sensitive content in edge justifications.

\paragraph{Ethical considerations.}
This work studies defensive interventions against documented attack patterns in a synthetic benchmark; no real user data or live systems are involved. The defense experiences are derived from sanitized benchmark trajectories, and the redaction pipeline removes credential-shaped strings before any LLM call. We publish model roles and prompts so that the community can audit judge dependence. The technique itself---distilling failure traces into advisory memory---is defensive in intent, but the same mechanism could in principle memorize and replay attacker behavior; we consider it important that the experience store is filtered, scoped, and auditable, and we flag misuse of failure-derived memory as a risk worth governing.
